\section{Metodología de Desarrollo}

\subsubsection{Justificación de la Metodología}

El desarrollo del presente sistema se rige por la metodología \textbf{Scrum},
marco ágil definido por Schwaber y Sutherland (2020) como un conjunto liviano de
reglas que ayuda a las personas, equipos y organizaciones a generar valor a través
de soluciones adaptativas para problemas complejos \cite{scrum_guide_2020}. La elecci\'on de Scrum se
fundamenta en tres criterios principales:

\begin{itemize}[leftmargin=1.5em]
  \item \textbf{Complejidad e incertidumbre del dominio.} Un sistema restaurantero
        omnicanal con sincronizaci\'on en tiempo real, m\'ultiples roles, flujos
        concurrentes de liquidaci\'on e integraci\'on con un motor anal\'itico distribuido
        (\textit{Apache Spark}) requiere ciclos cortos de retroalimentaci\'on para detectar
        inconsistencias en los requisitos antes de que se propaguen al c\'odigo. Scrum
        formaliza este ciclo a trav\'es de Sprints de duraci\'on fija (dos semanas en este
        proyecto).
  \item \textbf{Equipo reducido y multidisciplinario.} El equipo est\'a formado por
        dos integrantes con perfiles complementarios: uno orientado al an\'alisis,
        dise\~no y documentaci\'on (Natalia Anaya Ram\'irez), y otro al desarrollo e
        infraestructura tecnol\'ogica (Diego Villagr\'an Salazar). Scrum est\'a
        optimizado para equipos de tres a nueve personas y promueve la
        colaboraci\'on directa sin capas burocr\'aticas intermedias \cite{scrum_guide_2020}.
  \item \textbf{Entregables incrementales verificables.} Los m\'odulos del sistema
        (OE-01 a OE-13) son independientes en su mayor parte y pueden demostrarse
        en ambiente de staging al final de cada sprint, facilitando la validaci\'on
        temprana con el director del trabajo terminal.
\end{itemize}

\subsubsection{Adaptación del Marco Scrum al Proyecto}

La Tabla~\ref{tab:scrum_roles} describe la adaptaci\'on de los roles de Scrum al
contexto del presente trabajo terminal, considerando que el equipo de dos personas
asume responsabilidades compartidas en varios roles.

\begin{table}[ht]
\caption{Adaptación de Roles Scrum al Equipo del Proyecto}
\label{tab:scrum_roles}
\centering
\small
\setlength{\tabcolsep}{4pt}
\begin{tabular}{>{\bfseries\raggedright\arraybackslash}p{3cm}
                >{\raggedright\arraybackslash}p{4cm}
                >{\raggedright\arraybackslash}p{9cm}}
\hline
\thead{Rol Scrum} & \thead{Responsable} & \thead{Responsabilidades en el proyecto} \\
\hline
Product Owner &
  Natalia Anaya Ram\'irez &
  Define y prioriza el Product Backlog. Valida que los criterios de aceptaci\'on
  (evidencias de los OE) se cumplan al cierre de cada Sprint. \\ \rowcolor{altrow}
Scrum Master &
  Diego Villagr\'an Salazar &
  Facilita las ceremonias (Daily, Sprint Review, Retrospectiva). Elimina
  impedimentos t\'ecnicos (configuraci\'on de entornos, integraciones). \\ \rowcolor{altrow}
Development Team &
  Ambos integrantes &
  Natalia lidera an\'alisis, diagramas, documentaci\'on y dise\~no UX.
  Diego lidera arquitectura, desarrollo del stack y despliegue en Vercel/Supabase. \\ \rowcolor{altrow}
Director / Stakeholder &
  Guzm\'an Flores Jessie Paulina &
  Revisa avances en Sprint Reviews quincenales. Valida alineaci\'on con los
  objetivos del trabajo terminal. \\
\hline
\end{tabular}
\end{table}

Los \textbf{artefactos} principales son: (1)~\textit{Product Backlog}, formado
por los 13 objetivos espec\'ificos (OE-01 a OE-13) descompuestos en historias de
usuario; (2)~\textit{Sprint Backlog}, subconjunto de \'items seleccionados para
cada Sprint de dos semanas; y (3)~\textit{Incremento}, un entregable funcional
al final de cada Sprint desplegado en el ambiente de Preview de Vercel.

\subsubsection{Planificación por Sprints}

El desarrollo se estructura en \textbf{12 sprints} de duración fija (2--3 semanas cada uno), siguiendo la lógica incremental de Scrum. Cada sprint tiene un objetivo claro y entregables verificables:

\begin{enumerate}[leftmargin=*, itemsep=1pt, topsep=2pt]
  \item \textbf{Sprint 1:} Analítica de negocio y Arquitectura.
  \item \textbf{Sprint 2:} Diseño de la base de datos (MER y Modelo Lógico).
  \item \textbf{Sprint 3:} Desarrollo de la base de datos y migraciones iniciales.
  \item \textbf{Sprint 4:} Diseño de vistas, mockups y prototipado UX.
  \item \textbf{Sprint 5:} Diseño y arquitectura de backend (Server Actions y API).
  \item \textbf{Sprint 6:} Definición de KPI y vistas de Dashboard administrativo.
  \item \textbf{Sprint 7:} Implementaci\'on del pipeline de datos con Apache Spark (Arquitectura Medallion).
  \item \textbf{Sprint 8:} Desarrollo del m\'odulo para el personal operativo (Mesero).
  \item \textbf{Sprint 9:} Desarrollo del m\'odulo para el consumidor (Comensal).
  \item \textbf{Sprint 10:} Desarrollo del módulo para cocina (Tablero Chef).
  \item \textbf{Sprint 11:} Integración de la lógica de liquidación y control de pagos.
  \item \textbf{Sprint 12:} Pruebas integrales, ajustes de seguridad y cierre.
\end{enumerate}

Se utiliza un tablero Kanban en GitHub para la gestión visual de tareas, permitiendo el seguimiento en tiempo real del estado de cada funcionalidad (\textit{To Do}, \textit{In Progress}, \textit{Review}, \textit{Done}).

\subsubsection{Cronograma General (Diagrama de Gantt)}

El proyecto abarca del 1 de enero de 2026 al 30 de noviembre de 2026 y se divide
en dos grandes fases: (1)~\textbf{Fase Te\'orica y de Dise\~no} (enero--mayo), a
cargo de Natalia Anaya Ram\'irez, que cubre la documentaci\'on, el an\'alisis de
reglas de negocio, el dise\~no de la base de datos, los flujos, los diagramas UML,
la factibilidad y la arquitectura; y (2)~\textbf{Fase de Implementaci\'on y Pruebas}
(mayo--noviembre), a cargo de Diego Villagr\'an Salazar, que comprende la
configuraci\'on de infraestructura, el desarrollo de m\'odulos, la integraci\'on
y las pruebas.

La Tabla~\ref{tab:gantt} presenta el cronograma detallado con las actividades,
los responsables y la distribuci\'on temporal.

\begin{table}[H]
\caption{Cronograma del Proyecto --- Diagrama de Gantt (Enero--Noviembre 2026)}
\label{tab:gantt}
\centering
\scriptsize
\setlength{\tabcolsep}{1pt}
\begin{tabular}{>{\raggedright\arraybackslash}p{7.5cm}
                >{\centering\arraybackslash}p{0.5cm}
                >{\centering\arraybackslash}p{0.5cm}
                >{\centering\arraybackslash}p{0.5cm}
                >{\centering\arraybackslash}p{0.5cm}
                >{\centering\arraybackslash}p{0.5cm}
                >{\centering\arraybackslash}p{0.5cm}
                >{\centering\arraybackslash}p{0.5cm}
                >{\centering\arraybackslash}p{0.5cm}
                >{\centering\arraybackslash}p{0.5cm}
                >{\centering\arraybackslash}p{0.5cm}
                >{\centering\arraybackslash}p{0.5cm}
                >{\raggedright\arraybackslash}p{2cm}}
\hline
\textbf{Actividad} &
\rotatebox{90}{\textbf{Ene}} &
\rotatebox{90}{\textbf{Feb}} &
\rotatebox{90}{\textbf{Mar}} &
\rotatebox{90}{\textbf{Abr}} &
\rotatebox{90}{\textbf{May}} &
\rotatebox{90}{\textbf{Jun}} &
\rotatebox{90}{\textbf{Jul}} &
\rotatebox{90}{\textbf{Ago}} &
\rotatebox{90}{\textbf{Sep}} &
\rotatebox{90}{\textbf{Oct}} &
\rotatebox{90}{\textbf{Nov}} &
\textbf{Resp.} \\
\hline
\multicolumn{13}{l}{\textit{\textbf{Fase 1 — Teórica y de Diseño (Natalia Anaya Ramírez)}}} \\
\hline
\rowcolor{lightblue}
Planteamiento del problema y justificaci\'on &
  \cellcolor{accent}\textcolor{white}{\textbf{--}} & \cellcolor{accent}\textcolor{white}{\textbf{--}} &&&&&&&&&& Natalia \\
Revisi\'on de literatura y estado del arte &
  \cellcolor{accent}\textcolor{white}{\textbf{--}} & \cellcolor{accent}\textcolor{white}{\textbf{--}} &&&&&&&&&& Natalia \\
\rowcolor{lightblue}
Definici\'on de objetivos, alcances y l\'imites &
  \cellcolor{accent}\textcolor{white}{\textbf{--}} & \cellcolor{accent}\textcolor{white}{\textbf{--}} &&&&&&&&&& Natalia \\
Estudio de factibilidad t\'ecnica y operativa &
  & \cellcolor{accent}\textcolor{white}{\textbf{--}} & \cellcolor{accent}\textcolor{white}{\textbf{--}} &&&&&&&&& Natalia \\
\rowcolor{lightblue}
Estudio de factibilidad legal y econ\'omica &
  & \cellcolor{accent}\textcolor{white}{\textbf{--}} & \cellcolor{accent}\textcolor{white}{\textbf{--}} &&&&&&&&& Natalia \\
Matriz de riesgos y mitigaci\'on &
  && \cellcolor{accent}\textcolor{white}{\textbf{--}} & \cellcolor{accent}\textcolor{white}{\textbf{--}} &&&&&&&& Natalia \\
\rowcolor{lightblue}
Definici\'on de reglas de negocio &
  && \cellcolor{accent}\textcolor{white}{\textbf{--}} & \cellcolor{accent}\textcolor{white}{\textbf{--}} &&&&&&&& Natalia \\
Modelado de flujos de usuario y casos de uso &
  &&& \cellcolor{accent}\textcolor{white}{\textbf{--}} & \cellcolor{accent}\textcolor{white}{\textbf{--}} &&&&&&& Natalia \\
\rowcolor{lightblue}
Dise\~no de m\'aquinas de estado &
  &&& \cellcolor{accent}\textcolor{white}{\textbf{--}} & \cellcolor{accent}\textcolor{white}{\textbf{--}} &&&&&&& Natalia \\
Definici\'on de requerimientos (RF y RNF) &
  &&& \cellcolor{accent}\textcolor{white}{\textbf{--}} & \cellcolor{accent}\textcolor{white}{\textbf{--}} &&&&&&& Natalia \\
\rowcolor{lightblue}
Dise\~no de base de datos (MER y modelo l\'ogico) &
  &&&& \cellcolor{accent}\textcolor{white}{\textbf{--}} & \cellcolor{accent}\textcolor{white}{\textbf{--}} &&&&&& Natalia \\
Dise\~no de arquitectura del sistema &
  &&&& \cellcolor{accent}\textcolor{white}{\textbf{--}} & \cellcolor{accent}\textcolor{white}{\textbf{--}} &&&&&& Natalia/Diego \\
\rowcolor{lightblue}
Diagramas UML (clases, secuencia, despliegue, casos de uso) &
  &&&& \cellcolor{accent}\textcolor{white}{\textbf{--}} & \cellcolor{accent}\textcolor{white}{\textbf{--}} &&&&&& Natalia \\
Dise\~no de interfaces (wireframes / prototipos) &
  &&&& \cellcolor{accent}\textcolor{white}{\textbf{--}} & \cellcolor{accent}\textcolor{white}{\textbf{--}} &&&&&& Natalia \\
\rowcolor{lightblue}
\textbf{HITO 1: Entrega y presentaci\'on TT1} &
  &&&& \cellcolor{red}\textcolor{white}{\textbf{$\star$}} &&&&&&& \textbf{Ambos} \\
\hline
\multicolumn{13}{l}{\textit{\textbf{Fase 2 — Implementación y Pruebas (Diego Villagrán Salazar)}}} \\
\hline
\rowcolor{lightblue}
Configuraci\'on de infraestructura (Vercel, Supabase, CI/CD) &
  &&&&& \cellcolor{primary}\textcolor{white}{\textbf{--}} & \cellcolor{primary}\textcolor{white}{\textbf{--}} &&&&& Diego \\
Implementaci\'on schema Prisma y migraciones &
  &&&&& \cellcolor{primary}\textcolor{white}{\textbf{--}} & \cellcolor{primary}\textcolor{white}{\textbf{--}} &&&&& Diego \\
\rowcolor{lightblue}
Desarrollo m\'odulo autenticaci\'on y roles (OE-03) &
  &&&&&& \cellcolor{primary}\textcolor{white}{\textbf{--}} & \cellcolor{primary}\textcolor{white}{\textbf{--}} &&&& Diego \\
Desarrollo m\'odulo estructura org. y usuarios (OE-02, OE-04) &
  &&&&&& \cellcolor{primary}\textcolor{white}{\textbf{--}} & \cellcolor{primary}\textcolor{white}{\textbf{--}} &&&& Diego \\
\rowcolor{lightblue}
Desarrollo m\'odulo cat\'alogo y men\'u (OE-05) &
  &&&&&&& \cellcolor{primary}\textcolor{white}{\textbf{--}} & \cellcolor{primary}\textcolor{white}{\textbf{--}} &&& Diego \\
Desarrollo m\'odulo mesas f\'isicas y virtuales (OE-06, OE-07) &
  &&&&&&& \cellcolor{primary}\textcolor{white}{\textbf{--}} & \cellcolor{primary}\textcolor{white}{\textbf{--}} &&& Diego \\
\rowcolor{lightblue}
Desarrollo m\'odulo \'ordenes y canasta (OE-08) &
  &&&&&&&& \cellcolor{primary}\textcolor{white}{\textbf{--}} & \cellcolor{primary}\textcolor{white}{\textbf{--}} && Diego \\
Desarrollo tablero Chef en tiempo real (OE-09) &
  &&&&&&&& \cellcolor{primary}\textcolor{white}{\textbf{--}} & \cellcolor{primary}\textcolor{white}{\textbf{--}} && Diego \\
\rowcolor{lightblue}
Desarrollo m\'odulo liquidaci\'on y control concurrencia (OE-10) &
  &&&&&&&&& \cellcolor{primary}\textcolor{white}{\textbf{--}} & \cellcolor{primary}\textcolor{white}{\textbf{--}} & Diego \\
\rowcolor{lightblue}
Desarrollo del pipeline anal\'itico en Spark (Medallion) &
  &&&&&&&&& \cellcolor{primary}\textcolor{white}{\textbf{--}} & \cellcolor{primary}\textcolor{white}{\textbf{--}} & Diego \\
Desarrollo dashboard analytics (OE-12) &
  &&&&&&&&&& \cellcolor{primary}\textcolor{white}{\textbf{--}} & Diego \\
\rowcolor{lightblue}
Implementaci\'on medidas de seguridad y RLS (OE-13) &
  &&&&&&&&&& \cellcolor{primary}\textcolor{white}{\textbf{--}} & Diego \\
Pruebas unitarias e integraci\'on &
  &&&&&& \cellcolor{primary}\textcolor{white}{\textbf{--}} & \cellcolor{primary}\textcolor{white}{\textbf{--}} & \cellcolor{primary}\textcolor{white}{\textbf{--}} & \cellcolor{primary}\textcolor{white}{\textbf{--}} && Diego \\
\rowcolor{lightblue}
Pruebas de carga y rendimiento (k6) &
  &&&&&&&&& \cellcolor{primary}\textcolor{white}{\textbf{--}} & \cellcolor{primary}\textcolor{white}{\textbf{--}} & Diego \\
Correcciones y ajustes post-pruebas &
  &&&&&&&&&& \cellcolor{primary}\textcolor{white}{\textbf{--}} & Diego \\
\rowcolor{lightblue}
Documentaci\'on t\'ecnica final (TT2) &
  &&&&&&&&& \cellcolor{primary}\textcolor{white}{\textbf{--}} & \cellcolor{primary}\textcolor{white}{\textbf{--}} & Ambos \\
\textbf{HITO 2: Entrega y presentaci\'on TT2} &
  &&&&&&&&&& \cellcolor{red}\textcolor{white}{\textbf{$\star$}} & \textbf{Ambos} \\
\hline
\multicolumn{13}{p{15.5cm}}{\footnotesize\textit{$\star$ = Hito de entrega/presentaci\'on. Azul oscuro = Diego Villagr\'an (Fase~2). Azul claro = Natalia Anaya (Fase~1). TT1: mayo 2026. TT2: noviembre 2026.}} \\
\end{tabular}
\end{table}

% -----------------------------------------------------------
